\bigskip
\begin{center}
{\scriptsize
\renewcommand{\arraystretch}{1.1}
\setlength{\tabcolsep}{4pt}
\begin{tabularx}{\textwidth}{c l X c}

\multicolumn{4}{c}{\normalsize\textbf{小组成员信息与贡献占比}} \\[0.4em]

\toprule
\textbf{照片} & \textbf{姓名（学号）} & \textbf{主要贡献} & \textbf{占比} \\
\midrule

\includegraphics[width=1.3cm,height=1.6cm,keepaspectratio]{郭凡渊}
& 郭凡渊（202421215132）
& \raggedright
-- 前端界面开发与优化（主题切换、像素动画、样式美化）\\
-- 后端三层架构重构、代码结构优化与视觉模型集成\\
-- Live2D 角色资产生成与配置、指导文档体系建立\\
-- PPT 制作与答辩准备
& 26\% \\
\midrule

\includegraphics[width=1.3cm,height=1.6cm,keepaspectratio]{刘广裕}
& 刘广裕（202421215123）
& \raggedright
-- 论文撰写、参考文献整理与格式修订\\
-- API 接口规范测试、模拟后端开发与前端性能优化\\
-- 人设文档与表达映射编写、PPT 制作与答辩准备
& 24\% \\
\midrule

\includegraphics[width=1.3cm,height=1.6cm,keepaspectratio]{杨文博}
& 杨文博（202421215115）
& \raggedright
-- 论文初稿撰写与 LaTeX 排版\\
-- 项目安全访问控制与 LLM 推理优化\\
-- 配置常量管理与代码重构、PPT 制作与答辩准备
& 25\% \\
\midrule

\includegraphics[width=1.3cm,height=1.6cm,keepaspectratio]{温嘉浩}
& 温嘉浩（202421215121）
& \raggedright
-- 后端 Agent 工作流架构设计与 LangGraph 图编排\\
-- 后端文件布局、文档整理与 MiniMax 兼容性优化\\
-- 仓库维护与代码清理、PPT 制作与答辩准备
& 25\% \\
\midrule
\multicolumn{4}{p{\textwidth}}{
\scriptsize\raggedright\textit{注：所有成员贡献占比经小组协商一致确定，合计为 100\%。论文编写由杨文博（初稿与排版）和刘广裕（文献与修订）共同负责，郭凡渊与温嘉浩参与内容审阅。}
} \\
\bottomrule
\end{tabularx}
}
\end{center}
